\section{Terrain classification}\label{sec:terrain}

What kind of ground does a flight fly over, and how is this written for a whole flight?

\paragraph{The rule.} The class of a flight is not the altitude of its takeoff. Take every point of the flight on the
\UniGridS-s grid, in every segment of the cleaned track, from the takeoff to the landing. The ground elevation under a point
is read from the IGN BD ALTI v2 DEM of 25 m (Lambert-93), averaged over blocks of $8\times8$ cells on a grid of 200 m aligned
on multiples of 200 m: a point takes the mean of the block that holds it. A block with more than half of its cells without
data takes no value, and neither does a block of which all cells are exactly 0 m (the sea, which the DEM delivers as 0 m): a
point on such a block, or outside the DEM, is dropped. The tiles of the DEM are read one at a time and added into the blocks,
never the whole of France at 25 m. With $z_1,\dots,z_n$ the ground elevations of the points that remain, the
\emph{median ground elevation} of the flight is $\tilde z=\operatorname{median}(z_k)$ and its class is
\begin{align*}
  \text{plains}&:\ \tilde z<\UniEdgePlains\ \text{m}, &
  \text{hills}&:\ \UniEdgePlains\le\tilde z<\UniEdgeHills\ \text{m},\\
  \text{low mountains}&:\ \UniEdgeHills\le\tilde z<\UniEdgeLow\ \text{m}, &
  \text{high mountains}&:\ \tilde z\ge\UniEdgeLow\ \text{m}.
\end{align*}
A flight with less than \UniMinOnDem{} of its points on the DEM is \emph{unknown}: it is counted, kept in the raw comparison
(all terrain together) and left out of the comparison by class and by region. Fig.~\ref{fig:rule} shows the rule on one
flight.

\begin{figure}
\centering
\includegraphics[width=\textwidth]{terrain_rule.pdf}
\caption{The class of a flight is the class of the median ground elevation under its \UniGridS-s points, whatever its
takeoff altitude. A real paraglider flight: the ground under its points (grey), its altitude (blue), the median (black) and the
class edges (dashed).}\label{fig:rule}
\end{figure}

\paragraph{What it gives.} Table~\ref{tab:counts} counts the flights and the pilots of each cell. Of the \UniFunnelParaOpen{}
open paraglider flights \UniCellParaUnknownFlights{} are unknown, and \UniCellHangUnknownFlights{} of the \UniFunnelHangOpen{} hang glider flights.
Fig.~\ref{fig:hist} gives the median ground elevation of the flights with the edges; the median of the medians is
\UniMedianPara{} m for paragliders and \UniMedianHang{} m for hang gliders. Fig.~\ref{fig:map} places the takeoffs.

\begin{table}
\centering
\small
\renewcommand{\arraystretch}{1.15}\setlength{\tabcolsep}{4pt}
\begin{adjustbox}{max width=\textwidth}
\begin{tabular}{lrrrrrr}
\toprule
Discipline & Plains & Hills & Low mountains & High mountains & Unknown & All \\
\midrule
para & 12\,212 (1\,383) & 10\,414 (1\,930) & 21\,090 (4\,097) & 8\,867 (2\,713) & 3\,296 (747) & 55\,879 (5\,540) \\
hang & 174 (35) & 104 (43) & 782 (152) & 225 (81) & 132 (36) & 1\,417 (191) \\
\bottomrule
\end{tabular}
\end{adjustbox}
\caption{Flights (pilots) per discipline and terrain class. A pilot who flew in two classes counts in both; in the last column, once.}\label{tab:counts}
\end{table}


\begin{figure}
\centering
\includegraphics[width=\textwidth]{terrain_hist.pdf}
\caption{Hang gliders concentrate on low mountains, paragliders spread over every class. Median ground elevation of the open
flights that have one, per discipline, with the class edges (dashed).}\label{fig:hist}
\end{figure}

\begin{figure}
\centering
\includegraphics[width=\textwidth]{map.pdf}
\caption{The takeoffs concentrate in the Alps for both disciplines, and the hang glider takeoffs elsewhere are scattered. Takeoffs of
the open flights, coloured by the class of the flight, and the four regions of Section~\ref{sec:regions}.}\label{fig:map}
\end{figure}

\paragraph{Limits.} The DEM covers France: a flight that leaves it, over Switzerland or Italy, loses those points, and is
unknown when it spends more than half of its points there. The 200 m blocks smooth the ridges and the sea is not ground.

\keep{A flight is classified by the median ground elevation under its points, with edges $\UniEdgePlains$, $\UniEdgeHills$ and
$\UniEdgeLow$ m; flights with less than \UniMinOnDem{} of their points on the DEM are unknown and only enter the raw comparison.}
